% The local response package constructed at the selected scale.
\subsection{Construction of the finite local response package}
\label{sec:ex-lattice-trace-eve-producer}

We give the quantitative construction behind the local response package used
by the later trace estimates. The output contains its probe error, chart
error, bound, and cloud allowance; these are built together because the final
cloud billing uses the same bound budget as the probe estimate.

Fix the scale data of the trace construction and write
\[
 \chi(x)=\sqrt{H(x)},\qquad e=\mathrm{title}(P_1),\quad
 a=\mathrm{title}(P_0),\quad
 h=\mathrm{credibility}(P_1),\quad \nu=\mathrm{DoorsRate}(L).
\]
Here $H$ and $\chi$ are the monitor and its square root from
Subsection~\ref{sec:ex-lattice-trace-local-rate-selection}, whose functions
$\beta(x)$, $b(x)$, $R(x)$ and $\rho(x)$ are also used below; we write
$b(\beta(x))$ for the radius function evaluated at $\beta(x)$,
\[
 \operatorname{fieldCut}(x)=x^{-1/2}\exp\bigl((\log\log x)^2\bigr),
 \qquad
 \operatorname{fineCut}(x)=\operatorname{fieldCut}(\beta(x))
                        =\beta(x)^{-1/2}H(x).
\]
Further, $P_0$ and $P_1$
are the two stage profiles of the transition, and $L$ is the block side;
$e,a,h,\nu$ are positive, $e,a\ge1$, and $\nu=L^{-2/5}$. Let $\mathscr F$ be the grade-$d$ coefficient frame and $F$ its
finite family of coefficient-recovery probes; $F$ is distinct from the
finite basis of the frame, whose indices
are denoted by $\alpha$ below. The bounds $\mathcal B_{\rm test}$,
$\mathcal B_{\rm read}$, $\mathcal B_{\rm err}$ and $\mathrm{orbitCost}$ below
are constants of the frame $\mathscr F$ (the second and third also of the
selected holonomy map), not of the probe family. For $p\in F$, let $w_p=\pi(p)\ge1$ be its
probe bound, and write $\Lambda=M_d$ for the polarization constant
in the probe estimate, both as in
Subsection~\ref{sec:ex-lattice-trace-price-recurrence}. Put
\[
 Z=\mathcal B_{\rm test}(\mathscr F),\qquad T=\mathcal B_{\rm read}(\mathscr F),\qquad
 c=\mathrm{orbitCost}(\mathscr F)\,\mathcal B_{\rm err}(\mathscr F)\,\mathrm{corn}(d)\,a.
\]
The coefficient-recovery and readout constants $Z,T$ are chosen at
least one: the defining finite-dimensional inequalities remain valid
when their constants are enlarged to one. The error bound is the test
bound times one plus a finite sum of nonnegative chain-rule bounds;
the orbit cost is one plus a finite sum of coefficient sizes; and
$\mathrm{corn}(d)$ is one plus a finite sum of nonnegative comparison
constants. All three are at least one. Together with $a\ge1$ this gives
$c\ge1$, in particular $c>0$. These are elementary properties of the
chosen frame constants, not additional analytic inputs. For a positive target
bound $\zeta$, define
\begin{equation}
 \kappa=\min\left\{\frac{\nu}{2cZ},\zeta\right\}>0,
 \qquad
 g_p(s)=w_p\Lambda L^d\,2(es+s)+s.
 \label{eq:ex-lattice-eve-gp}
\end{equation}
Finiteness of the probe family $F$ is essential here: each $g_p$ is continuous,
$g_p(0)=0$, and therefore there is one $b>0$ such that $g_p(s)\le\kappa$
for all $p\in F$ and $0\le s\le b$. Set $Q=b/2$.

\begin{theorem}[Finite-probe and cloud billing]
\label{thm:ex-lattice-eve-producer}
Let $d\in\mathbb N$, and assume $e,a,h,\nu,c,Z,T,\zeta>0$, $e,a,Z,T\ge1$,
and a finite probe family $F$ with weights $w_p\ge0$ and normalization
$\Lambda\ge0$. Suppose the selected response satisfies
$u\in\mathcal R_{-1/2}$ and $\chi\in\mathcal R_0$, with $\chi(x)\ge1$
eventually.
For each probe, the charged scalar is the continuous function
\eqref{eq:ex-lattice-eve-gp}. Set
\[
 \kappa=\min\left\{\frac{\nu}{2cZ},\zeta\right\},\qquad
 \delta_x=Q/\chi(x),\qquad C_x=\kappa/\chi(x).
\]
Assume the following selected-scale receipts hold on one common eventual
tail in $x$; the tail is fixed before the probe, the frame index and all
the other data listed in each receipt are chosen, and each receipt holds
uniformly in those data. For the threshold $\theta=\operatorname{fieldCut}(x)$
and the fine threshold $\Theta=\operatorname{fineCut}(x)$,
\begin{equation}
 \tfrac12\le \frac{\Theta}{\theta}\le2.
 \label{eq:ex-lattice-eve-blend}
\end{equation}
The accompanying spatial comparison is
$b(\beta(x))\le2b(x)$ on this same tail. Indeed, that section gives $\beta(x)=x+c_L$ for a fixed $c_L$ and
$b(x)=\lceil(\log x)^2\rceil$. Thus
$(\log(x+c_L))^2/(\log x)^2\to1$, while rounding changes either
quantity by at most one. The ratio $b(\beta(x))/b(x)$ tends to one,
and is eventually at most two.
Both comparisons are used when applying the local budget estimate:
the fine threshold is within a factor two of the coarse threshold,
and the fine budget radius fits the allowed coarse radius.
The probe receipt is as follows. Let $p\in F$ have order $r_p$ and probe
functional $R_p$ on coefficients, and let $\alpha$ be a basis index of the
coefficient frame, with basis coefficient $\mathsf e_\alpha$. Consider any
step whose block count is $L$ and whose first-order row is literally the
row of the selected kernel with block side $L$, with coarse side larger
than $2\rho(x)$, and any two thresholds $w,y\in(0,1]$ within a factor two
of each other. Let $U_\alpha$ be the coarse-lattice function obtained from the
local presentation of $\mathsf e_\alpha$ at the origin with threshold $w$
by cropping the coarse configuration to radius $\rho(x)$ about the parent
of the origin and routing it through the first-order row of the step; let
$V_p$ be the logarithmic-link test curve of $p$ with amplitude $y$; and
let $\mathcal L$ be the leading coefficient map determined by the selected
kernel and the two thresholds. The
\emph{probe-to-run difference} is
\[
 L^d\left.\frac{d^{r_p}}{ds^{r_p}}U_\alpha(V_p(s))\right|_{s=0}
   -R_p\bigl(\mathcal L\,\mathsf e_\alpha\bigr),
\]
and the receipt is that its absolute value is at most $Q/\chi(x)$ for
every $p\in F$, every $\alpha$, and all such steps and thresholds. The cap
receipt is: for every support $X$ and every pair of budgets $(P,B)$ such
that $P$ has threshold $\theta$ and the power of $P_0$, $B$ has threshold
$\Theta$ and radius $b(\beta(x))$, and
$\operatorname{coverNumber}(B.B,X)\le1$, and for every
$\beta\ge\beta(x)/2$, the scale cap $O_x$ (the response cost of
\eqref{eq:ex-lattice-explicit-local-cost} without its last, paint, line)
and the local moment $M_x$ (the
head-count moment of order $4(n+1)$ there) obey
\begin{equation}
 0\le O_x\le1,\qquad
 \mathrm{scaleRate}(e)\,O_x\,M_x\le \frac{Q}{\chi(x)}.
 \label{eq:ex-lattice-eve-cap}
\end{equation}
For every cloud $Z_0$ of mass $m_0\le h$, put
$\mathcal M_x(Z_0)=\mathrm{marvel}_x(Z_0)\ge0$. Its cloud receipt is
\begin{equation}
 L^d\,2\mathcal M_x(Z_0)\le S m_0,
 \qquad S=\frac{\nu}{2cTh}.
 \label{eq:ex-lattice-eve-cloud}
\end{equation}
Then for all sufficiently large $x$, with the same uniformity in the data
of each receipt:

\begin{enumerate}
\item the two selected thresholds satisfy the factor-two comparison;
\item the response amplitude $s_x=|u(x)|$ obeys $s_x\le\delta_x$ and
      $s_x\le\zeta/\chi(x)$;
\item the scale cap is at most one, and its clock-weighted moment is at most
      $\delta_x$;
\item every selected probe has approximation error at most $\delta_x$;
\item for each probe $p$, the explicit charge
      $w_p\Lambda L^d\,2(e s_x+\delta_x)+\delta_x$ is at most $C_x$;
\item for every cloud $Z_0$ of mass at most $h$, the combined charge
\[
 c\left[Z C_x+T L^d\,2\mathcal M_x(Z_0)\right]\le\nu.
\]
\end{enumerate}
In particular $C_x\le\zeta/\chi(x)$.
\end{theorem}

\begin{proof}
The product rate gives $\chi u\in\mathcal R_{-1/2}$. Choose
$q\in(-1/2,0)$; the strict-power definition implies
$\chi(x)|u(x)|\le\beta(x)^q\to0$. Thus, eventually,
\begin{equation}
 |u(x)|\le Q/\chi(x),\qquad |u(x)|\le\zeta/\chi(x),\qquad \chi(x)\ge1.
 \label{eq:ex-lattice-eve-response-bound}
\end{equation}
The first two give the asserted response amplitude bound and the desired
estimate for the amplitude. The first selected geometry estimate gives
the factor-two threshold comparison.

The selected scale-cap estimate is called at bound $Q/e$. Its estimate
is linear in the allowed scalar multiplier $t$; evaluating it at $t=e$
gives exactly the second inequality in \eqref{eq:ex-lattice-eve-cap}, while its companion bound gives
$O_x\le1$. The probe approximation estimate is called at tolerance $Q$,
which gives the bound stated above directly.

For a probe $p$, the amplitude bound
\eqref{eq:ex-lattice-eve-response-bound} gives
\begin{equation}\label{eq:ex-lattice-eve-amplitude}
 2(e|u(x)|+\delta_x)
       \le \frac{2(eQ+Q)}{\chi(x)}.
\end{equation}
Consequently its charged expression is bounded by
\[
 \frac{w_p\Lambda L^d\,2(eQ+Q)+Q}{\chi(x)}
   =\frac{g_p(Q)}{\chi(x)}
   \le \frac{\kappa}{\chi(x)}=C_x.
\]
This is the simultaneous finite-probe choice made above; it is not a
pointwise choice depending on the probe.

Finally, let a cloud have mass $m_0\le h$. The cloud estimate \eqref{eq:ex-lattice-eve-cloud} gives
\[
 T L^d(2\mathcal M_x(Z_0))\le T S h=\frac{\nu}{2c}.
\]
The probe part costs at most $ZC_x\le Z\kappa/\chi(x)\le Z\kappa$.
By the definition of $\kappa$, $cZ\kappa\le\nu/2$. Adding the cloud and
probe charges therefore gives at most $\nu$. This proves all six clauses.
\end{proof}

\subsubsection{Matching the finite-volume response estimate}
\label{subsec:ex-lattice-eve-matching}

We now identify the local insertion receipt in the actual response
construction. Fix a leaf insertion of weight $\tau$, and hold its masked
ingredient fixed while varying the local field. The two functions compared
are the actual-sector and Gaussian-reference integrals
\begin{equation}
 \Phi_6(y)=\int R_6(y,z)\chi_6(y,z)h(z)\,d\gamma_{C_0}(z),\qquad
 \Phi_9(y)=\int R_9(y,z)\chi_9(y,z)h(z)\,d\gamma_{C_0}(z),
 \label{eq:ex-lattice-eve-two-integrals}
\end{equation}
where $h$ is the fixed measurable masked insertion and $|h|\le\tau$.
The density comparison in Subsection~\ref{sec:principal-gaussian-comparison}
proves, under its displayed finite-dimensional covariance, mean, cutoff,
and chart hypotheses,
\begin{equation}
 \|\Phi_6-\Phi_9\|_{J^n_0}
 \le \tau P_n(N,Q,S)\bigl(\delta+2N e^{-v+M/2}\bigr).
 \label{eq:ex-lattice-eve-linear-receipt}
\end{equation}
The primitive mean and covariance estimates supplying $\delta$ are proved
in Subsection~\ref{sec:principal-comparison-jets}. The fixed insertion is not
differentiated. The local chart directions are the same ordered variations
used in the response test, with $j\le n$ and their prescribed directional
and canopy bounds.

In the exact construction, the fixed leaf is held together with its regular
masked ingredient $L'$. For each fixed grade and block, the bounds
$N(x),Q(x),S(x),\delta(x)$ and the cutoff-tail bounds are common to every
admissible period, forest and tag choice, every leaf whose cover is at most
$B$, and every normalized direction list allowed in the comparison. They
are independent of the leaf weight $\tau$. Thus the one finite-probe choice
and selected response $u(x)$ work uniformly over all these inputs. The
local read set, forest, tags, insertion mask, and chart are fixed during each
individual differentiation. The two local observables are
\begin{equation}
 \mathcal O_{\mathrm{act}}(y)=\Phi_6(y),\qquad
 \mathcal O_{\mathrm{ref}}(y)=\Phi_9(y)
 \label{eq:ex-lattice-eve-source-identification}
\end{equation}
and these identities follow by writing each side with its own normalized
Gaussian density and using the protected-support agreement on the
insertion's read set. Thus the
response difference is exactly $\Phi_6-\Phi_9$, not a comparison
of different insertions or different Gaussian spaces. The local finite
bound cutoff makes the integrands vanish outside a common compact coordinate
ball; the cutoffs are one on the protected smaller ball. Differentiation
under the integral is justified with the bounded measurable $h$ fixed, by
the uniform derivative bounds on that compact ball and dominated
convergence.

The selected rate budgets choose the scalar envelope
\[
 B_*(x)=P_n(N(x),Q(x),S(x))
             \bigl(\delta(x)+2N(x)e^{-v(x)+M(x)/2}\bigr).
\]
The covariance discrepancy has strict-power rate $-1/2$, the mean
discrepancy has rate $-1$ before conversion to Gaussian coordinates, and
the good-coordinate tail has every negative rate. The fixed polynomial
factor has rate zero. Therefore $B_*\in\mathcal R_{-1/2}$.
The response construction packages the concrete derivative, cutoff, and
tail budgets into one scalar envelope. The component bounds are assembled
from the finite product and sum rules for the displayed rates. With this
envelope, for every admissible leaf and every ordered list of at most $n$
normalized chart directions, the derivative of the actual-minus-reference
observable at the chart origin is bounded by $\tau B_*(x)$, uniformly over
the choices above. This is the concrete uniform derivative inequality used
by the response construction.

We can now choose the scalar in the insertion receipt explicitly. The
polynomial majorant $P_n$ has nonnegative coefficients, and its arguments,
$\delta$ and the cutoff-tail term are nonnegative. Thus $B_*\ge0$ on
the common tail where the comparison is applied. Define
\begin{equation}\label{eq:ex-lattice-eve-concrete-amplitude}
 f(x)=\tfrac12 B_*(x),\qquad u(x)=2|f(x)|=B_*(x)
\end{equation}
on that tail, and set both functions to zero before it. Altering a
function on an initial interval does not change its eventual rate.
Multiplication by a fixed positive constant preserves the strict-power
rate: for each $q>-1/2$, choose $q'\in(-1/2,q)$ and absorb the
constant using $\beta(x)^{q-q'}\to\infty$. Hence both $f$ and $u$
belong to $\mathcal R_{-1/2}$.

For every admissible leaf of weight $\tau\ge0$ and every ordered list
of at most $n$ allowed chart directions, the inequality already proved
is therefore exactly
\[
 \bigl\|D^{\mathbf v}(\mathcal O_{\mathrm{act}}
                  -\mathcal O_{\mathrm{ref}})(0)\bigr\|
 \le\tau B_*(x)=\tau\,2|f(x)|=\tau|u(x)|.
\]
This includes the zeroth derivative and the case $\tau=0$; no division
by the leaf weight is made. The density and cutoff derivative bounds
used in the comparison also give the required differentiability near
the chart origin, with the masked ingredient fixed. Thus the bound
applies to the actual response test, not just to a formal coefficient
of its Taylor expansion.

The function $f$ may depend on the fixed grade, block, stage profiles,
selected kernel, differentiation order and comparison parameters. It is
chosen before the period, forest, tags, leaf, mask, read support, base
configuration and allowed direction list. Its tail is uniform in all
those latter data, because the envelope $B_*$ and the two integral
identities were uniform in them. There is no additional normalization,
positivity or matching condition on the auxiliary scalar $f$: it is a
bound majorant, not one of the local fields. The numerical factor two
is therefore realized by the definition in
\eqref{eq:ex-lattice-eve-concrete-amplitude}; it incurs no extra
analytic loss.

If no branch patch compatible with the fixed selected kernel and block
exists, take $f=u=0$ everywhere. Any step to which the receipt could
apply would itself supply such a patch, so the universally quantified
insertion receipt has no instances in this case. If a compatible patch
exists, the actual step uses that selected patch and the preceding
construction applies. In either case this supplies the insertion
receipt. The finite-test choice and scalar billing of
Theorem~\ref{thm:ex-lattice-eve-producer} now apply to this concrete
$u$, on the intersection of its common tail with the geometry, cap,
probe and cloud tails; no outstanding insertion estimate is assumed.

The package for one transition with block side $L$ and grade $d$ is
constructed under the following standing conditions, which the scalar
theorem above does not see: the guard of the record satisfies
$D.\varepsilon\le1/4$; the two integer parameters of the selected kernel
are at least $4$, and the first of them is at least $d$; the differentiation
order satisfies $n\ge2$; and the real parameter $\mu$ of the response
estimate is positive. The response receipt with amplitude $|u(x)|$ is stated there for
every leaf whose label is the logarithmic key of a regular activity and
whose ingredient is exactly the mask, at the threshold of the
representation, of a function that is regular on the support of that
activity; for every budget of radius $b(x)$ and of the power of $P_0$ that
is compatible with the representation; for every read support that is
healthy at the base configuration and has cover number at most that
radius; and for every list of at most $n$ directions satisfying the
budget's directional bound and the canopy bound
$2\Theta(1+R(x))^{p}$, where $p$ is the power of $P_0$. These identities
of the leaf, mask, budget and support are what make
\eqref{eq:ex-lattice-eve-source-identification} a comparison of one
insertion on one Gaussian space. The cap receipt of
Theorem~\ref{thm:ex-lattice-eve-producer} follows from the scalar budget
estimate \eqref{eq:ex-lattice-budget-output} of
Subsection~\ref{sec:ex-lattice-trace-local-rate-selection}, called
at bound $Q/e$ and with the nonnegative paint term omitted. In fact
$B.\mathrm{threshold}=\Theta=\operatorname{fineCut}(x)$ and
$P.\mathrm{threshold}=\theta\le2\Theta$ by \eqref{eq:ex-lattice-eve-blend};
$B.B=b(\beta(x))\le2b(x)$ by the comparison proved above;
$P.\mathrm{power}$ is the power of $P_0$; and the cap's cover bound
is at most one, so the support allowance in
\eqref{eq:ex-lattice-budget-support} can be chosen to be the constant one.
The lower coupling bound is the same in both statements. These verify
all the support and budget conditions of that application. The two remaining receipts are proved next for these same local data.

\subsubsection{The actual probe and cloud bounds}
\label{subsec:trace-eve-probe-cloud}
Fix a probe $p$ of order $k$ and a frame column $\alpha$ of
amplitude degree $\ell_\alpha$. The constants below depend on these
finite choices and the installed kernel and block, but not on the
period or on the two thresholds $w,y\in(0,1]$ with
$1/2\le y/w\le2$. Let $R=\rho(x)=45(1+\lceil(\log x)^2\rceil)$.
We apply the cropped-probe calculation of
Subsection~\ref{sec:primary-calibration-producer} with fine threshold
$w$ and coarse probe threshold $y$. Its actual curve is exactly the
one defining $U_\alpha(V_p(s))$: it crops to radius $R$ and applies
the installed first-order row. The room condition $2R<c$ ensures
that the crop uses the canonical covering-lattice lifts. The
full polynomial curve uses the same row. Absolute convergence of
its Fourier kernel identifies the periodized convolution with the
covering sum regrouped by residue classes.

In the notation of \eqref{eq:primary-probe-tail}, the jets of the
two normalized input curves through order $k$ differ by at most
\[
 e_R=C\{1+K(1+4R)^{dk}\}\exp(-\sigma R/2),\qquad \sigma>0.
\]
Here the polynomial factor bounds the probe's link-log jets outside
the crop; inside the crop the two arrays agree. The constants are
uniform in $w,y$: each nonempty chain-rule block contributes a
threshold factor $y^r/w\le y/w\le2$.
The full predicted jets have a fixed bound $U$ given by the weighted
kernel moments. The normalized local-section estimate
\eqref{eq:primary-stencil-stability} therefore bounds the output
jet difference, including the physical factor, by
\begin{equation}\label{eq:trace-eve-probe-envelope}
 E_{p,\alpha}(x)
   =L^d C_{k,\alpha,U}(1+e_R)^{\ell_\alpha+k}e_R.
\end{equation}
The full polynomial forecast in that section has exactly the
coefficient transport $\mathcal L\mathsf e_\alpha$: its three
fine-grade factors multiply to $L^d w^{-\ell_\alpha}$, and its
coarse scaling inserts $y$ in the probe. Thus its reading is
$R_p(\mathcal L\mathsf e_\alpha)$, with no change of coefficient
map. This identifies the difference estimated in
\eqref{eq:trace-eve-probe-envelope} with the probe-to-run difference
in Theorem~\ref{thm:ex-lattice-eve-producer}. For the unit column the
local section is constant and both readings agree exactly.

Since $R\asymp(\log x)^2$, the logarithm of $e_R$ is at most
$-c_0(\log x)^2+O(\log\log x)$ for some $c_0>0$.
Also $\log\chi(x)=\tfrac12(\log\log\beta(x))^2$ and
$\beta(x)=x+c_L$. It follows directly that
$\chi(x)E_{p,\alpha}(x)\to0$. There are finitely many probes
and columns. Taking the maximum of their eventual thresholds gives
$E_{p,\alpha}(x)\le Q/\chi(x)$ simultaneously, uniformly over
all the steps, periods and thresholds allowed in the receipt.
This proves the probe estimate with its required order of quantifiers.

For the cloud, let $q_*<0$ be the fixed rate of the selected target
and $\epsilon_*>0$ its support reserve. The actual cloud response
is the quantity in Lemma~\ref{lem:ex-lattice-medal-cloud-price}, with
its coefficient now specified explicitly:
\begin{equation}\label{eq:trace-eve-cloud-formula}
 \mathcal M_x(Z_0)
 =m_0\exp(-\epsilon_*N(Z_0))\,\beta(x)^{q_*/4},
 \qquad N(Z_0)\ge0.
\end{equation}
The coefficient is a negative power of $\beta$, so it has the
strict-power rate required in that lemma. In particular
$0\le\mathcal M_x(Z_0)\le m_0\beta(x)^{q_*/4}$ uniformly
over clouds and representations. Choose one scale threshold with
$2L^d\beta(x)^{q_*/4}\le S=\nu/(2cTh)$.
Multiplying by $m_0\ge0$ proves \eqref{eq:ex-lattice-eve-cloud}, including the case $m_0=0$
without division. The same formula gives the stronger statement
\[
 \mathcal M_x(Z_0)\le (v/\chi(x))m_0
\]
for every prescribed $v>0$ on a later common tail, since
$\chi(x)\beta(x)^{q_*/4}\to0$. This is the cloud estimate
used for regular and primary atoms. Thus the probe, cap, cloud,
and insertion estimates are all supplied before their common
finite intersection is used to construct the response package.
